\documentclass[10pt,a4paper]{amsart}
\usepackage[margin=19mm]{geometry}
\usepackage{amsmath,amssymb,mathtools}
\usepackage[scr=rsfs,cal=euler]{mathalfa}
\usepackage{xfrac}
\usepackage[unicode,hidelinks,bookmarks=false]{hyperref}
\newcommand{\ie}{i.e.\ }
\newcommand{\cf}{cf.\ }
\newcommand{\resp}{respectively\ }
\newcommand{\Subsection}{\subsection}
\providecommand{\nobreakdash}{\nobreak\hbox{-}}
\setlength{\emergencystretch}{2em}
\multlinegap0pt
\usepackage{mathtools}
\usepackage{amssymb}
\usepackage{bm}
\usepackage{enumerate}
\usepackage[shortlabels]{enumitem}
\usepackage{tikz}
\usepackage{stackengine}
\newtheorem{lemma}{Lemma}
\newtheorem{prop}{Proposition}
\newcommand{\bS}{\mathbf{S}}
\newcommand{\calM}{\mathcal{M}}
\newcommand{\id}{\operatorname{id}}
\newcommand{\ob}{\operatorname{ob}}
\title{Galois actions on surfaces and a higher genus Grothendieck-Teichm\"{u}ller group}
\author{Luciana Basualdo Bonatto, Marcy Robertson}
\date{Source: arXiv:2606.01466v1 (CC BY 4.0)}
\begin{document}
\maketitle

\noindent\emph{Source and licence.} Galois actions on surfaces and a higher genus Grothendieck-Teichm\"{u}ller group. Version arXiv:2606.01466v1 (CC BY 4.0), distributed by arXiv under the Creative Commons Attribution 4.0 International licence, http://creativecommons.org/licenses/by/4.0/, as recorded in the arXiv licence metadata for this exact version and shown on the abstract page https://arxiv.org/abs/2606.01466v1. Full licence notice: \texttt{licenses/arxiv-2606.01466-CC-BY-4.0.txt}.\par\smallskip

\noindent\textit{Editorial scope.} Counted unit: the article's prop prop:S-generated-by-standard-moves (source0.tex lines 2216--2218). The article's complete original proof of it is delivered with it (source0.tex lines 2220--2296, one \texttt{proof} environment of 77 lines). No mathematics is written, repaired or re-derived: the statement and the proof are the article's own lines, in the article's own notation, and no retained line is substituted or re-flowed.  Retained uncounted as the source-local context the proof needs: lines 2142--2144.  Nothing was omitted from any retained span, so the package carries no omission record.  No retained line contains a citation, so the wrapper carries no bibliography and no bibliography entry.  No retained line contains a figure environment, an image-inclusion command or drawing code, so no image is delivered and the package declares no asset dependency.  Editorial formatting only, in addition: an AMS \texttt{amsart} preamble that restates the article's own declarations and macros used by the retained lines, namely the environments \texttt{lemma}, \texttt{prop} on separate counters, together with the standard packages those lines need.  The retained environments print in this document as Lemma 1, Proposition 1 in order of appearance, whereas the article prints the same items with the numbering of its own full document.  The complete native source of the article as distributed in the arXiv e-print for this exact version is bundled unchanged as \texttt{sources/201828/source0.tex}.
\begin{lemma}\label{lemma: the complexes of morphisms and markings}
Fix a graph $G\in \ob(\bS(g,n+1))$ and let $S_G$ be the associated labelled, parametrized surface. Then the coslice category $G\downarrow \bS(g,n+1)$ is isomorphic to the groupoid $\mathbb{M}_G$.
\end{lemma}

\begin{prop}\label{prop:S-generated-by-standard-moves}
Every morphism in $\bS$ can be obtained from $\tau$, $\beta$, $\alpha$, $\sigma$, and their inverses using categorical composition, operadic composition, symmetric group actions, and contractions.
\end{prop}

\begin{proof} 
Let $[\phi]:G\to H$ be a morphism in $\bS(g,n+1)$, represented by a boundary-label-preserving diffeomorphism
\[
\phi:S_G\longrightarrow S_H.
\]
By Lemma~\ref{lemma: the complexes of morphisms and markings}, this morphism corresponds to the marked pants decomposition of $S_G$ obtained by pulling back the standard marked pants decomposition of $S_H$ along $\phi$.

We first translate a single edge of the complex of markings into a morphism of $\bS$. Suppose that
\[
(C,[m])\rightsquigarrow (C',[m'])
\]
is an edge of $\calM(S_G)$ of type $T$, $B$, $A$, or $S$. Let
\[
[\phi]:G\to H
\qquad\text{and}\qquad
[\psi]:G\to H'
\]
be the objects of the coslice category $G\downarrow\bS(g,n+1)$ corresponding to $(C,[m])$ and $(C',[m'])$, respectively, under Lemma~\ref{lemma: the complexes of morphisms and markings}. The edge determines the unique morphism in the coslice category from $[\phi]$ to $[\psi]$; equivalently, there is a unique morphism $[f]:H\to H'$ such that
\[
[\psi]=[f]\circ[\phi].
\]
We claim that $[f]$ is obtained from identities and one of the standard generators $\tau,\beta,\alpha,\sigma$ by operadic composition, symmetric group actions, and contractions.


We verify the claim for a $B$-move; the other cases are analogous. Suppose that
\[
(C,[m]:H\hookrightarrow S_G)
\overset{B_{a,b,c}}{\rightsquigarrow}
(C',[m']:H'\hookrightarrow S_G)
\]
is a $B$-move. Then $C=C'$, so the two marked pants decompositions have the same underlying cut system. The move is supported on a single cut-subsurface
\[
S'\subset \overline{S_G\setminus C}
\]
with boundary components $a,b,c$, and the markings $m$ and $m'$ agree on the complement of $S'$.

Choose an identification $\xi:S'\xrightarrow{\cong}P$ with the standard pair of pants such that
\[
\xi\circ m|_{S'}=m_P,\qquad
\xi(a)=\partial_1,\qquad
\xi(b)=\partial_2,\qquad
\xi(c)=\partial_0.
\]
By the definition of the $B$-move, under this identification the new marking is given by
\[
\xi\circ m'|_{S'}=\beta\circ \xi\circ m|_{S'}.
\]

Let $v\in V(H)$ and $v'\in V(H')$ be the vertices corresponding to the component $S'$. Since the move is supported on $S'$, the graphs $H$ and $H'$ agree away from these vertices, with all boundary labels preserved. We therefore define a mapping class
\[
f:S_H\longrightarrow S_{H'}
\]
by taking it to be the identity on every pair-of-pants component except the one corresponding to $v$, and by taking it on that component to be the standard $B$-move $\beta$ under the above identifications. If the local identification with the standard pair of pants differs by a cyclic relabelling, then the local map is the corresponding cyclic translate of $\beta$.

Thus $[f]:H\to H'$ is obtained from $\beta$ by operadic composition with identities, together with the relevant symmetric group action. Moreover, the marking corresponding to $[\psi]$ is exactly the marking obtained from the one corresponding to $[\phi]$ by applying $[f]$ after $[\phi]$. Hence
\[
[\psi]=[f]\circ[\phi].
\]


We now complete the proof. Let $[\phi]\in \bS(g,n+1)(G,H)$, and let $(C_\phi,[m_\phi])$ be the corresponding marked pants decomposition of $S_G$. Since $\calM(S_G)$ is connected, there is an edge path
\[
(C_0,[m_0])\rightsquigarrow (C_1,[m_1])
\rightsquigarrow \cdots \rightsquigarrow (C_k,[m_k])
\]
from the standard marking $(C_0,[m_0])=(iE_G,[m_G])$ to $(C_k,[m_k])=(C_\phi,[m_\phi])$.

Under Lemma~\ref{lemma: the complexes of morphisms and markings}, the standard marking corresponds to the identity morphism $\id_G:G\to G$. By the single-edge argument above, each edge in the path corresponds to postcomposition with a morphism
\[
[f_i]:H_{i-1}\to H_i
\]
obtained from identities and one of $\tau,\beta,\alpha,\sigma$ by operadic composition, symmetric group actions, and contractions. Hence
\[
[\phi]=[f_k]\circ\cdots\circ [f_1]\circ \id_G.
\]
Thus $[\phi]$ lies in the subcategory generated by $\tau,\beta,\alpha,\sigma$ under categorical composition, operadic composition, symmetric group actions, and contractions. Since $[\phi]$ was arbitrary, the claim follows.
\end{proof}

\end{document}
